% Part: first-order-logic
% Chapter: beyond
% Section: many-sorted-logic

\documentclass[../../../include/open-logic-section]{subfiles}

\begin{document}

\olfileid{fol}{byd}{msl}

\olsection{المنطق متعدد الأصناف}

في منطق الرتبة الأولى تتراوح المتغيرات والمكممات على !!{domain}
واحد. وقد يقتضي الغرض تعدد ال!!{domain}s وتباينها؛ فنريد مثلًا
!!a{domain} للأعداد، و!!a{domain} للأشياء الهندسية،
و!!a{domain} للدوال بين الأعداد، و!!a{domain} للزمر الإبدالية،
وغير ذلك.

ولهذا يوضع المنطق متعدد الأصناف. نبدأ بقائمة «أصناف»، فصنف
الشيء يعيّن ال!!{domain} الذي ينتمي إليه. ثم نجعل لكل صنف
!!{variable}s ومكممات، ونجعل له في المعتاد !!{identity} أيضًا.
ونحدّد أنواع الدوال والعلاقات بحسب أصناف حججها الجائزة.
وفيما عدا ذلك تبقى قواعد منطق الرتبة الأولى على حالها، مع
إفراد كل صنف بصور قواعد المكممات.

فلو أردنا مثلًا دراسة العلاقات الدولية، جاز أن نختار لغة ذات
صنفين: المواطنين الفرنسيين والمواطنين الألمان. ثم نضع علاقة
أحادية «يشرب النبيذ» للصنف الأول، وأحادية أخرى «يأكل النقانق»
للثاني، وثنائية «يؤلفان زوجًا متعدد الجنسية» حجتها الأولى من
الصنف الأول والثانية من الثاني. فإذا جعلنا المتغيرات $a$ و$b$
و$c$ للفرنسيين، والمتغيرات $x$ و$y$ و$z$ للألمان، كانت الصيغة
\[
\lforall[a][\lforall[x][(\Atom{\Obj{MarriedTo}}{a,x} \lif
(\Atom{\Obj{DrinksWine}}{a} \lor \lnot \Atom{\Obj{EatsWurst}}{x})]]
\]
تقرر أن من تزوّج من الفرنسيين شخصًا ألمانيًا، فإما أن الفرنسي
يشرب النبيذ، وإما أن الشخص الألماني لا يأكل النقانق.

ويُضمَّن المنطق متعدد الأصناف في منطق الرتبة الأولى بجمع أشياء
ال!!{domain}s المختلفة في !!{domain} واحد، ثم تمييز الأصناف
ب!!{predicate}s أحادية وتقييد المكممات بها. ففي اللغة المقابلة
للمثال المتقدم نزيد ال!!{predicate}s الأحادية
«$\Obj{German}$» و«$\Obj{French}$»، ونبقي العلاقات الأخرى
بعد رفع قيود الأصناف من تركيبها. ويُترجم المكمم المصنف
$\lforall[x][!A]$، حيث $x$ !!a{variable} ألماني الصنف، إلى
\[
\lforall[x][(\Atom{\Obj{German}}{x} \lif !A)].
\]
ثم نزيد بديهيات تمنع اجتماع الصنفين، مثل
$\lforall[x][\lnot (\Atom{\Obj{German}}{x} \land
  \Atom{\Obj{French}}{x})]$، وبديهيات تقضي بأن «يشرب النبيذ»
لا يصدق إلا على ما يحقق المحمول $\Atom{\Obj{French}}{x}$،
وعلى ذلك يُقاس. وبهذه الاصطلاحات والبديهيات يسهل بيان أن
ال!!{sentence}s المتعددة الأصناف تُترجم إلى !!{sentence}s من
الرتبة الأولى، وأن ال!!{derivation}s المتعددة الأصناف تُترجم إلى
!!{derivation}s من الرتبة الأولى. وكذلك تقابل ال!!{structure}s
المتعددة الأصناف !!{structure}s من الرتبة الأولى؛ ويصح الرجوع
من الثانية إلى الأولى إذا قسمت محمولات الأصناف المجال إلى
أصناف غير خالية متباينة، واستوفت الرموز قيود الأنواع.
ومن هذا التقابل نحصل على مبرهنة التمام للمنطق متعدد الأصناف.

\end{document}
